%!TEX root = 00_main.tex
% ============================================================================
% ADVANCED SLIDES - CODE, VIDEOS, QR CODES, GANTT CHARTS
% ============================================================================

\section{Code Listings}

% --- Python Code ---
\begin{frame}[fragile]{Python Code Example}
    \begin{lstlisting}[language=Python, caption={Wind power calculation}]
import numpy as np

def wind_power(rho, A, Cp, v):
    """Calculate wind turbine power output.

    Args:
        rho: Air density (kg/m^3)
        A: Rotor area (m^2)
        Cp: Power coefficient
        v: Wind speed (m/s)
    """
    return 0.5 * rho * A * Cp * v**3

# Example: 5 MW turbine
P = wind_power(1.225, 12469, 0.45, 12)
print(f"Power: {P/1e6:.2f} MW")
    \end{lstlisting}
\end{frame}

% --- MATLAB Code ---
\begin{frame}[fragile]{MATLAB Code Example}
    \begin{lstlisting}[language=Matlab, caption={Blade element momentum theory}]
function [a, ap] = bem_iteration(lambda, theta, Cl, Cd)
    % BEM iteration for axial and tangential induction
    sigma = 0.05;  % Solidity
    phi = atan(1 / lambda);

    % Thrust and torque coefficients
    Ct = sigma * (Cl*cos(phi) + Cd*sin(phi));
    Cq = sigma * (Cl*sin(phi) - Cd*cos(phi));

    % Induction factors
    a = 1 / (1 + 4*sin(phi)^2 / (sigma*Cl*cos(phi)));
    ap = 1 / (4*cos(phi) / (sigma*Cl) - 1);
end
    \end{lstlisting}
\end{frame}

% --- Inline Code ---
\begin{frame}[fragile]{Code Snippets and Terminal Output}
    \textbf{Inline code:} Use \lstinline|numpy.fft.fft()| for spectral analysis.

    \vspace{0.5cm}

    \textbf{Terminal output:}
    \begin{lstlisting}[style=terminal]
$ python simulate.py --turbine IEA15MW --wind 12
Loading turbine model... done
Running simulation: 600s at dt=0.01s
  Progress: [====================] 100%
Results saved to: results/IEA15MW_12ms.h5
Peak thrust: 2.45 MN
DEL (flapwise): 14.2 MNm
    \end{lstlisting}
\end{frame}

\section{Multimedia}

% --- Video Placeholder ---
\begin{frame}[fragile]{Video Embedding}
    \centering

    % Placeholder for video (uncomment and adjust path when you have a video)
    % \movie[width=0.8\textwidth, height=0.5\textwidth, poster, showcontrols]
    %     {\includegraphics[width=0.8\textwidth]{example_figure.png}}
    %     {animation.mp4}

    \includegraphics[width=0.6\textwidth]{example_figure.png}

    \vspace{0.3cm}

    {\small\textit{Video placeholder -- Replace with:}}

    \begin{lstlisting}[basicstyle=\tiny\ttfamily, frame=none]
\movie[width=0.8\textwidth, poster, showcontrols]
    {\includegraphics[width=0.8\textwidth]{poster.png}}
    {video.mp4}
    \end{lstlisting}

    {\small Requires PDF viewer with multimedia support (Adobe Reader).}
\end{frame}

% --- Animation with Overlays ---
\begin{frame}{Animation with Overlays}
    \centering
    \textbf{Blade rotation simulation}

    \vspace{0.5cm}

    \begin{tikzpicture}
        % Hub
        \fill[gray] (0,0) circle (0.3);

        % Blades at different angles (simulated rotation)
        \only<1>{\draw[line width=3pt, mainblue] (0,0) -- (90:2.5);}
        \only<1>{\draw[line width=3pt, mainblue] (0,0) -- (210:2.5);}
        \only<1>{\draw[line width=3pt, mainblue] (0,0) -- (330:2.5);}

        \only<2>{\draw[line width=3pt, mainblue] (0,0) -- (120:2.5);}
        \only<2>{\draw[line width=3pt, mainblue] (0,0) -- (240:2.5);}
        \only<2>{\draw[line width=3pt, mainblue] (0,0) -- (0:2.5);}

        \only<3>{\draw[line width=3pt, mainblue] (0,0) -- (150:2.5);}
        \only<3>{\draw[line width=3pt, mainblue] (0,0) -- (270:2.5);}
        \only<3>{\draw[line width=3pt, mainblue] (0,0) -- (30:2.5);}

        \only<4>{\draw[line width=3pt, mainblue] (0,0) -- (180:2.5);}
        \only<4>{\draw[line width=3pt, mainblue] (0,0) -- (300:2.5);}
        \only<4>{\draw[line width=3pt, mainblue] (0,0) -- (60:2.5);}
    \end{tikzpicture}

    \vspace{0.3cm}
    Frame \only<1>{1}\only<2>{2}\only<3>{3}\only<4>{4} of 4
\end{frame}

\section{QR Codes}

% --- QR Code Examples ---
\begin{frame}{QR Codes for References}
    \begin{columns}[T]
        \begin{column}{0.3\textwidth}
            \centering
            \qrcode[height=2.5cm]{https://www.nrel.gov/wind/}

            \vspace{0.2cm}
            {\small NREL Wind Research}
        \end{column}
        \begin{column}{0.3\textwidth}
            \centering
            \qrcode[height=2.5cm]{https://github.com/OpenFAST/openfast}

            \vspace{0.2cm}
            {\small OpenFAST GitHub}
        \end{column}
        \begin{column}{0.3\textwidth}
            \centering
            \qrcode[height=2.5cm]{https://doi.org/10.1016/j.renene.2020.01.001}

            \vspace{0.2cm}
            {\small Paper DOI}
        \end{column}
    \end{columns}

    \vspace{0.5cm}

    {\small Scan QR codes to access resources, papers, or code repositories.}
\end{frame}

% --- QR Code with Contact ---
\begin{frame}{Contact Information}
    \begin{columns}[T]
        \begin{column}{0.5\textwidth}
            \textbf{Contact Details:}

            \vspace{0.3cm}

            \begin{itemize}
                \item[] \textbf{Name:} \AuthorName
                \item[] \textbf{Institution:} \AuthorAffiliation
                \item[] \textbf{Email:} name@example.com
                \item[] \textbf{ORCID:} 0000-0000-0000-0000
            \end{itemize}

            \vspace{0.5cm}

            Scan for digital business card $\rightarrow$
        \end{column}
        \begin{column}{0.4\textwidth}
            \centering
            \qrcode[height=4cm]{https://orcid.org/0000-0000-0000-0000}
        \end{column}
    \end{columns}
\end{frame}

\section{Project Planning}

% --- Gantt Chart ---
\begin{frame}{Project Timeline (Gantt Chart)}
    \centering
    \begin{ganttchart}[
        hgrid,
        vgrid,
        x unit=0.6cm,
        y unit chart=0.5cm,
        title height=1,
        bar height=0.6,
        bar/.style={fill=mainblue},
        bar incomplete/.style={fill=mainblue!30},
        milestone/.style={fill=red, shape=diamond},
        title label font=\tiny,
        bar label font=\tiny,
        group label font=\tiny\bfseries
    ]{1}{12}
        \gantttitle{2026}{12} \\
        \gantttitlelist{1,...,12}{1} \\

        \ganttgroup{WP1: Design}{1}{4} \\
        \ganttbar{Conceptual}{1}{2} \\
        \ganttbar{Detailed}{2}{4} \\

        \ganttgroup{WP2: Simulation}{3}{8} \\
        \ganttbar{Model setup}{3}{4} \\
        \ganttbar{Load cases}{4}{7} \\
        \ganttbar{Post-processing}{7}{8} \\

        \ganttgroup{WP3: Validation}{8}{11} \\
        \ganttbar{Testing}{8}{10} \\
        \ganttbar{Certification}{10}{11} \\

        \ganttmilestone{Final Report}{12}

        \ganttlink{elem2}{elem3}
        \ganttlink{elem5}{elem6}
        \ganttlink{elem8}{elem9}
    \end{ganttchart}
\end{frame}

% --- Detailed Gantt ---
\begin{frame}{Detailed Work Package Schedule}
    \centering
    \begin{ganttchart}[
        hgrid,
        vgrid,
        x unit=0.45cm,
        y unit chart=0.4cm,
        title height=1,
        bar height=0.5,
        bar/.style={fill=mainblue},
        milestone/.style={fill=red, shape=diamond, inner sep=1pt},
        title label font=\tiny,
        bar label font=\tiny,
    ]{1}{16}
        \gantttitle{Q1}{4}
        \gantttitle{Q2}{4}
        \gantttitle{Q3}{4}
        \gantttitle{Q4}{4} \\
        \gantttitlelist{1,...,16}{1} \\

        \ganttbar{Literature review}{1}{3} \\
        \ganttbar{Method development}{2}{6} \\
        \ganttbar{Implementation}{5}{10} \\
        \ganttbar{Testing}{9}{13} \\
        \ganttbar{Documentation}{12}{15} \\
        \ganttmilestone{Submission}{16}

        \ganttlink{elem1}{elem2}
        \ganttlink{elem2}{elem3}
        \ganttlink{elem3}{elem4}
        \ganttlink{elem4}{elem5}
    \end{ganttchart}

    \vspace{0.3cm}
    {\small PhD project timeline -- Year 1}
\end{frame}
